\clearpage
\section{Kiến thức nền: đọc công thức ma trận như đọc một câu văn}
\label{sec:foundations}

\subsection{Vì sao cần thống kê nhiều chiều?}
Giả sử trên mỗi người ta đo huyết áp và nhịp tim. Mỗi người cung cấp một
\emph{cặp} số, không phải hai quan sát rời rạc. Người có nhịp tim cao có thể
cũng có huyết áp cao. Nếu phân tích từng biến riêng, ta bỏ qua mối liên hệ này.
Thống kê nhiều chiều giữ nguyên các thành phần thuộc cùng một đối tượng và
phân tích chúng đồng thời.

Một quan sát có $p$ biến được viết thành vectơ cột
\[
 \mathbf X=\begin{pmatrix}X_1\\X_2\\\vdots\\X_p\end{pmatrix}\in\R^p.
\]
Ở đây $p$ là \emph{số biến}. Khi lấy mẫu $n$ đối tượng, ta có
$\mathbf X_1,\ldots,\mathbf X_n$, mỗi vectơ vẫn chứa $p$ thành phần.
Chỉ số dưới của $\mathbf X_j$ chỉ \emph{đối tượng}; chỉ số dưới của $X_i$ trong
một vectơ chỉ \emph{biến}. Nếu cần cả hai, viết $X_{ji}$: biến thứ $i$ của đối
tượng thứ $j$.

\begin{example}[Đọc một bảng dữ liệu]
\[
 \mathbf D=\begin{pmatrix}2&12\\8&9\\6&9\\8&10\end{pmatrix}.
\]
Mỗi hàng là một đối tượng, mỗi cột là một biến: $n=4$, $p=2$.
Quan sát đầu là vectơ cột $\mathbf x_1=(2,12)'$.
Trong bài giảng, chữ hoa như $\mathbf X_1$ chỉ đại lượng ngẫu nhiên,
chữ thường như $\mathbf x_1$ chỉ giá trị đã quan sát.
\end{example}

\kiemtra{Một bảng có 60 hàng và 4 cột cho $n$ và $p$ bằng bao nhiêu?
Đáp án: $n=60$, $p=4$, không phải ngược lại.}

\subsection{Chuyển vị, phép nhân và kích thước}
Dấu phẩy trên $'$ hoặc ký hiệu $\mathsf T$ có nghĩa là \emph{chuyển vị}:
đổi hàng thành cột. Với $\mathbf a=(a_1,a_2)'$, ta có
$\mathbf a'=(a_1,a_2)$. Trong dữ liệu thực của bài giảng, hai ký hiệu này
có cùng nghĩa.
\[
 \mathbf a'\mathbf x=(a_1,a_2)
 \begin{pmatrix}x_1\\x_2\end{pmatrix}=a_1x_1+a_2x_2.
\]
Vậy $\mathbf a'\mathbf X$ chính là một tổng có trọng số. Ta có thể dùng nó
để biểu diễn tổng điểm, trung bình có trọng số hoặc hiệu giữa hai phép đo.

Nếu $A$ có kích thước $q\times p$, phép nhân $A\mathbf X$ có kích thước
$q\times1$. Quy tắc là: hai kích thước ở giữa phải khớp nhau.
\[
 \underbrace{A}_{q\times p}
 \underbrace{\mathbf X}_{p\times1}
 =\underbrace{\mathbf Y}_{q\times1}.
\]
Chẳng hạn,
\[
 A=\begin{pmatrix}1&-1\\1&1\end{pmatrix}
 \quad\Longrightarrow\quad
 A\mathbf X=\begin{pmatrix}X_1-X_2\\X_1+X_2\end{pmatrix}.
\]
Một phép nhân ma trận ở đây chỉ gói hai tổ hợp tuyến tính vào cùng một biểu thức.
Ma trận đơn vị $I_p$ giữ nguyên vectơ: $I_p\mathbf x=\mathbf x$.
Ma trận không $0$ cho kết quả bằng không, với kích thước phù hợp.

\luuy{Khi gặp một công thức khó, kiểm tra kích thước trước. Ví dụ
$\mathbf a'\Sigma\mathbf a$ có kích thước
$(1\times p)(p\times p)(p\times1)=1\times1$, nên là một số.}

\subsection{Nghịch đảo và định thức}
Nghịch đảo $A^{-1}$ là phép ``hoàn tác'' của $A$:
$A^{-1}A=AA^{-1}=I$. Chỉ ma trận vuông không suy biến mới có nghịch đảo.
Với ma trận hai chiều,
\[
 A=\begin{pmatrix}a&b\\c&d\end{pmatrix},\qquad
 |A|=ad-bc,
\]
\[
 A^{-1}=\frac1{ad-bc}\begin{pmatrix}d&-b\\-c&a\end{pmatrix}
 \quad\text{khi }ad-bc\ne0.
\]
Định thức $|A|$ là một số, không phải ma trận trị tuyệt đối của từng phần tử.
Nó đo hệ số thay đổi thể tích của phép biến đổi tuyến tính, xét theo trị tuyệt đối.
Nếu $|A|=0$, phép biến đổi làm mất ít nhất một hướng và không thể hoàn tác.

\begin{example}[Một nghịch đảo sẽ dùng nhiều lần]
\[
 \Sigma=\begin{pmatrix}4&3\\3&4\end{pmatrix},\qquad
 |\Sigma|=16-9=7,\qquad
 \Sigma^{-1}=\frac17\begin{pmatrix}4&-3\\-3&4\end{pmatrix}.
\]
Nhân hai ma trận sẽ được $I_2$. Khi tính bằng phần mềm, thay vì tạo nghịch đảo
một cách tường minh, thường giải hệ $\Sigma\mathbf z=\mathbf d$ rồi tính
$\mathbf d'\mathbf z$; đây là cùng đại lượng $\mathbf d'\Sigma^{-1}\mathbf d$.
\end{example}

\subsection{Kỳ vọng và ma trận hiệp phương sai}
Vectơ trung bình là
\[
 \vmu=\E(\mathbf X)=\begin{pmatrix}\E(X_1)\\\vdots\\\E(X_p)\end{pmatrix}.
\]
Ma trận hiệp phương sai là
\begin{equation}
 \Sigma=\Cov(\mathbf X)
 =\E\bigl[(\mathbf X-\vmu)(\mathbf X-\vmu)'\bigr].
 \label{eq:covdef}
\end{equation}
Tích bên trong là \emph{tích ngoài}: $(p\times1)(1\times p)=p\times p$.
Phần tử hàng $i$, cột $k$ là
\[
 \sigma_{ik}=\Cov(X_i,X_k)
 =\E[(X_i-\mu_i)(X_k-\mu_k)].
\]
Trên đường chéo, $\sigma_{ii}=\Var(X_i)$ là phương sai. Ngoài đường chéo là
hiệp phương sai; đối xứng vì $\sigma_{ik}=\sigma_{ki}$.
\nguon{Chương 4, §4.2, tr.~150, nhắc lại định nghĩa (2-30), (2-31) từ Chương 2}

Hiệp phương sai dương cho biết hai biến có xu hướng lệch cùng chiều so với
trung bình; hiệp phương sai âm cho biết chúng thường lệch ngược chiều.
Độ lớn phụ thuộc đơn vị đo, nên để so sánh người ta dùng tương quan
\[
 \rho_{ik}=\frac{\sigma_{ik}}{\sqrt{\sigma_{ii}\sigma_{kk}}},
 \qquad -1\le\rho_{ik}\le1,
\]
khi cả hai phương sai dương. Nếu hai biến độc lập và có moment bậc hai hữu hạn,
hiệp phương sai bằng không. Chiều ngược lại cần thêm giả thiết, sẽ được giải thích
ở Phần~\ref{sec:linear}.

\begin{example}[Đọc tham số]
Với $\vmu=(10,20)'$ và $\Sigma=\begin{pmatrix}4&3\\3&9\end{pmatrix}$,
biến thứ nhất có trung bình 10, độ lệch chuẩn 2; biến thứ hai có trung bình 20,
độ lệch chuẩn 3. Tương quan là $3/(2\cdot3)=0.5$.
Không đọc 3 ở ngoài đường chéo thành độ lệch chuẩn.
\end{example}

\subsection{Xác định dương: có độ biến thiên theo mọi hướng}
Ma trận đối xứng $\Sigma$ được gọi là \emph{xác định dương} nếu
\[
 \mathbf a'\Sigma\mathbf a>0\qquad\text{với mọi }\mathbf a\ne0.
\]
Vì $\mathbf a'\Sigma\mathbf a=\Var(\mathbf a'\mathbf X)$, điều này có nghĩa
mọi tổ hợp tuyến tính không tầm thường đều có phương sai dương. Không có một
hướng nào bị ``dẹt hoàn toàn''. Ma trận xác định dương có các trị riêng dương,
có nghịch đảo và định thức dương.

Một ma trận hiệp phương sai nói chung chỉ cần \emph{nửa xác định dương}:
$\mathbf a'\Sigma\mathbf a\ge0$. Chẳng hạn, nếu $X_2=2X_1$, thì
$2X_1-X_2$ luôn bằng 0; phương sai theo hướng $(2,-1)'$ bằng 0 và $\Sigma$
không khả nghịch. Trong bài giảng, trừ khi nói khác, các công thức có
$\Sigma^{-1}$ giả định $\Sigma$ xác định dương.

Với ma trận đối xứng $2\times2$, điều kiện xác định dương tương đương
$\sigma_{11}>0$ và $|\Sigma|>0$. Ở kích thước tổng quát, chỉ kiểm tra định
thức dương cho một ma trận đối xứng bất kỳ là chưa đủ; ví dụ $\diag(-1,-1)$
có định thức 1 nhưng không xác định dương. Riêng với ma trận \emph{đã biết là
hiệp phương sai}, định thức dương loại bỏ trường hợp suy biến.

\subsection{Ôn suy diễn một chiều: từ mẫu đến tổng thể}
Cho $X_1,\ldots,X_n$ iid, tức độc lập và cùng phân phối. Trung bình mẫu và
phương sai mẫu không chệch là
\[
 \bar X=\frac1n\sum_{j=1}^n X_j,\qquad
 s^2=\frac1{n-1}\sum_{j=1}^n(X_j-\bar X)^2.
\]
Độ lệch chuẩn $s$ mô tả độ phân tán giữa các quan sát; sai số chuẩn
$s/\sqrt n$ mô tả độ bất định của \emph{trung bình mẫu}. Hai đại lượng
phục vụ hai câu hỏi khác nhau.

Nếu mẫu lấy từ $\N(\mu,\sigma^2)$, dưới $H_0:\mu=\mu_0$,
\[
 t=\frac{\bar X-\mu_0}{s/\sqrt n}\sim t_{n-1}.
\]
Số bậc tự do $n-1$ phản ánh việc một mức trung bình đã được ước lượng từ mẫu.
Kiểm định hai phía mức $\alpha$ bác bỏ $H_0$ nếu
$|t|>t_{n-1;1-\alpha/2}$. Khoảng tin cậy tương ứng là
\[
 \bar x\pm t_{n-1;1-\alpha/2}\frac{s}{\sqrt n}.
\]
\nguon{Chương 5, §5.2, tr.~210--211, công thức (5-1)--(5-3)}

\paragraph{Quy ước phân vị xuyên suốt tài liệu.}
Nếu $Y$ có phân phối tích lũy $G$, thì $q_\gamma$ thỏa
$\Prob(Y\le q_\gamma)=\gamma$. Do đó $t_{\nu;0.975}$ là phân vị dưới
97.5\%, có xác suất đuôi phải 2.5\%. Tương tự,
$F_{r,s;0.95}$ và $\chi^2_{p;0.95}$ có xác suất đuôi phải 5\%.
Sách thường viết $F_{r,s}(\alpha)$ cho \emph{phân vị đuôi trên}
$\alpha$; nó tương ứng với $F_{r,s;1-\alpha}$ trong bài giảng.

\paragraph{Mức ý nghĩa, giá trị p và kết luận.}
Mức ý nghĩa $\alpha$ là xác suất bác bỏ sai giả thuyết không, khi giả thuyết
đó đúng và mô hình áp dụng đúng. Giá trị p là xác suất, \emph{dưới $H_0$},
nhận được thống kê ít nhất cực đoan như giá trị quan sát. Nó không phải
xác suất $H_0$ đúng. Nếu $\text{p-value}>\alpha$, viết ``chưa đủ bằng chứng
để bác bỏ'', không viết ``đã chứng minh $H_0$ đúng''.

Khoảng tin cậy 95\% có nghĩa: nếu lặp lại cách lấy mẫu và dựng khoảng,
khoảng 95\% các khoảng sẽ chứa tham số cố định. Với một khoảng đã tính,
không diễn giải 95\% thành xác suất hậu nghiệm của tham số khi chưa có mô hình Bayes.

\begin{example}[Một phép tính một chiều]
Với cột đầu $(2,8,6,8)'$ ở bảng dữ liệu trên,
$\bar x=6$ và $s^2=(16+4+0+4)/3=8$.
Kiểm định $\mu_0=7$ cho $t=(6-7)/\sqrt{8/4}=-0.7071$.
Khoảng 95\% là $6\pm3.1824\sqrt2$, tức xấp xỉ $[1.499;10.501]$.
Đây là suy diễn cho \emph{một} biến đã chọn, chưa phải phát biểu đồng thời
cho cả hai biến trong bảng.
\end{example}

\kiemtra{Nếu tăng $n$ lên bốn lần mà độ phân tán tổng thể giữ nguyên,
sai số chuẩn của trung bình giảm bao nhiêu? Đáp án: giảm một nửa.}
